% Submission version for Computer-Aided Design (Elsevier), built from the same section files as ../p1.tex.
\documentclass[preprint,12pt]{elsarticle}
\usepackage{amsmath,amssymb,amsthm,mathtools}
\usepackage{booktabs,array}
\usepackage{xcolor}
\usepackage[colorlinks=true,linkcolor=blue!50!black,citecolor=blue!50!black,urlcolor=blue!50!black]{hyperref}
\usepackage[margin=1in]{geometry}

% Drafting markers (none should remain in the submission).
\newcommand{\unverified}[1]{\textcolor{red}{[#1 --- not yet checked against the original]}}
\newcommand{\draftnote}[1]{\textcolor{orange!70!black}{[Draft note: #1]}}

\newtheorem{theorem}{Theorem}[section]
\newtheorem{lemma}[theorem]{Lemma}
\newtheorem{proposition}[theorem]{Proposition}
\newtheorem{corollary}[theorem]{Corollary}
\theoremstyle{definition}
\newtheorem{definition}[theorem]{Definition}
\newtheorem{example}[theorem]{Example}
\theoremstyle{remark}
\newtheorem{remark}[theorem]{Remark}

\newcommand{\R}{\mathbb{R}}
\newcommand{\Z}{\mathcal{Z}}
\newcommand{\C}{\mathcal{C}}
\newcommand{\wfs}{\operatorname{wfs}}
\newcommand{\dist}{\operatorname{dist}}
\newcommand{\inter}{\operatorname{int}}
\newcommand{\cl}{\operatorname{cl}}
\newcommand{\conv}{\operatorname{conv}}

\journal{Computer-Aided Design}
\emergencystretch=3em

\begin{document}
\begin{frontmatter}

\title{When does a tolerant B-rep define a solid? Existence, non-existence and uniqueness of realizations}

\author{Guoqing Zhang}
\address{ZENFLUX, Miami Gardens, Florida, USA}

\begin{abstract}
Boundary representations in CAD kernels carry tolerances: edge curves, the boundary curves of adjacent faces and vertex
points do not fit together exactly, and the solid they describe is only implied. We define a realization of a tolerant
B-rep as a PL embedding of its cell complex whose cells stay in their tolerance zones, and study existence,
non-existence and uniqueness. Uniqueness: if every shell has genus at least one, the carriers of the cells are
contractible, and each shell has a product collar containing its zones and a reference surface whose inclusion is a
homotopy equivalence, then the images of all realizations are topologically ambient isotopic, and so are the solids
they bound. The collar condition follows from a distance bound below the weak feature size of a PL reference surface,
for which we give an explicit lower bound in terms of conical radii, with the key lemma machine-checked in Lean.
Existence: a conditional fan-and-cone construction for regular complexes, and a re-intersection construction that
turns an embedded re-trimmed candidate into a realization. Non-existence: a checkable criterion for crossing planar
patches. In OpenCascade 8.0.1, the validity checks decide none of these questions: a sewn wedge passes both checks and
has no realization, and a thick pipe that passes both has two realizations whose solids are not homeomorphic.
\end{abstract}

\begin{keyword}
tolerant boundary representation \sep geometric robustness \sep realization \sep ambient isotopy \sep weak feature
size \sep 3-manifold topology
\end{keyword}

\end{frontmatter}

\section{Introduction}\label{sec:intro}

\subsection{The problem}
A solid in a CAD kernel is stored as a boundary representation: a cell complex (faces, edges, vertices) with
geometry attached to each cell. Each face carries a trimmed patch of a surface, each edge a 3D curve together
with one parameter-space curve (pcurve) for every adjacent face, and each vertex a point. In floating point,
and after modelling operations such as Booleans, fillets, offsets and sewing, these data do not agree: the
images of the pcurves on the surfaces miss the 3D edge curve, and the edge curves miss the vertex points. Kernels
in the tradition of Jackson~\cite{Jackson95} therefore attach a tolerance to every vertex, edge and face, and
regard the data as consistent when the deviations lie within the tolerances. OpenCascade, for instance, requires
the pcurve images to lie within the edge tolerance of the 3D curve, and the edge ends within the vertex tolerance
of the vertex point~\cite{OCCT}.

What solid does such a model describe? Qi and Shapiro~\cite{QiShapiro06b} describe the conventional answer
as: the representation is valid if there exists an exact solid in the tolerance zones. Two questions follow.
\begin{enumerate}
\item \emph{Existence.} Is there an embedded surface that respects all the tolerances?
\item \emph{Uniqueness.} If there is more than one, do they all bound the same solid up to ambient isotopy?
\end{enumerate}
Downstream algorithms (Booleans, offsets, meshing, analysis) implicitly assume both answers are yes.

\subsection{What we prove}
We call an embedding as in (1) a \emph{realization}, and make it precise in Section~\ref{sec:setting}: a PL
embedding $h$ of the cell complex $K$ with $h(\inter\sigma)\subseteq\Z(\sigma)$ for every cell $\sigma$, where
$\Z(\sigma)$ is the tolerance zone of $\sigma$.
\begin{itemize}
\item \textbf{Uniqueness} (Theorem~\ref{thm:uniqueness}). Suppose the carriers of the cells are contractible, each
shell has a product collar that contains its zones, and every shell has genus at least one. Then any two
realizations are ambient isotopic, and so are the solids they bound. The proof combines a carrier lemma (a
cellular homotopy argument) with Waldhausen's theorem that incompressible surfaces in a product $F\times I$ are
horizontal~\cite[Prop.~3.1]{Waldhausen68}.
\item \textbf{The collar condition is computable} (Section~\ref{sec:collar}). It holds when all zones lie within
distance $d$ of a reference surface $M$ with $d<\wfs(M)$, the weak feature size. For a polyhedral $M$ we prove
$\wfs(M)\ge\lambda(M)>0$, where $\lambda$ is defined from the conical radii of $M$ (Corollary~\ref{cor:lambda}). The
underlying cone lemma is machine-checked in Lean~4 with Mathlib. We also give a certified octree computation of
a lower bound for $\wfs(M)$.
\item \textbf{Existence} (Section~\ref{sec:existence}). There are two constructions, each under computable
hypotheses:
  \begin{itemize}
  \item fans in the normal disks of the edges, with cones at the vertices (Theorem~\ref{thm:lemma0});
  \item re-intersection of the face surfaces, for sharp edges with tight tolerances (Propositions~\ref{prop:B}
  and~\ref{prop:Bprime}).
  \end{itemize}
\item \textbf{Non-existence} (Theorem~\ref{thm:obstruction}). If two planar patches cross transversally away from
the projections (``shadows'') of all edge and vertex zones, there is no realization.
\end{itemize}

\subsection{What the experiments show}
In OpenCascade~8.0.1 (Section~\ref{sec:experiments}):
\begin{itemize}
\item \texttt{BRepCheck\_Analyzer} accepts models regardless of the size of the tolerances.
\item \texttt{BRepAlgoAPI\_Check} gives false alarms on realizable models, and misses models whose tolerances
exceed the weak feature size.
\item A thick pipe that passes both checks has two realizations whose solids are not homeomorphic: one ties a
trefoil into the hole.
\item Sewing a wedge whose faces are slightly displaced produces a model that passes both checks and, by
Theorem~\ref{thm:obstruction}, has no realization. The hypotheses were verified in exact rational arithmetic.
\end{itemize}

\subsection{Related work}
\begin{itemize}
\item \emph{Tolerances in practice.} Segal~\cite{Segal90} and Jackson~\cite{Jackson95} introduced vertex and
edge tolerances into polyhedral and B-rep modelling; Jackson's paper is the basis of the local-tolerance
approach used by commercial kernels.
\item \emph{Rigorous validity, without tolerances.} Sakkalis, Shen and Patrikalakis~\cite{SSP01} study interval
solids: boxes covering the boundary of a known exact solid, up to homeomorphism.
\item \emph{Global tolerances.} Qi and Shapiro~\cite{QiShapiro06} develop $\varepsilon$-solidity, a theory with
a single global tolerance that does not require an exact solid to exist~\cite{QiShapiro06b}. Shapiro~\cite{Shapiro08}
studies homotopy conditions for tolerant geometric queries.
\item \emph{Constructing an exact solid.} Andersson, Stewart and Zidani~\cite{ASZ07} and Stewart and
Zidani~\cite{SZ08} construct a solid from inconsistent trimmed-patch data by perturbing the patches (Whitney
extension, transfinite interpolation). They \emph{assume} that the perturbed patches form an embedded surface:
ASZ say that adjacent patches need an a posteriori check~\cite[pp.~819--820]{ASZ07}, and SZ make a similar
assumption~\cite[p.~188]{SZ08}.
\item \emph{Isotopy of approximations.} The weak feature size and its stability theory are due to Chazal and
Lieutier~\cite{ChazalLieutier07}; reach is due to Federer~\cite{Federer59}.
\end{itemize}
Not yet read: Andersson, Peters and Stewart~\cite{APS00} on equivalence of topological form; \unverified{the
2006 paper of Andersson, Stewart and Zidani}.

\emph{Novelty, stated cautiously.}
\begin{itemize}
\item Theorem~\ref{thm:uniqueness} applies Waldhausen's theorem; the new part is the set-up (zones, carriers,
collars) that makes it applicable to tolerant B-reps.
\item Positivity of $\wfs$ for compact semialgebraic sets is known: di Rocco, Eklund and
Gäfvert~\cite[Thm.~2.5]{dREG22} state it and attribute it to Fu~\cite{Fu85} \unverified{Fu's paper}. We did not find the explicit bound $\lambda(M)$ in the literature we checked. That does not make it
new.
\item The existence results derive the embedding from checkable hypotheses rather than assuming it. As far as
we have read, this is new; it remains to be checked against~\cite{APS00} and the 2006 ASZ paper.
\item We did not find an obstruction result like Theorem~\ref{thm:obstruction} in the literature we read.
\end{itemize}

\section{Tolerant B-reps and realizations}\label{sec:setting}

\subsection{Complexes}
$K$ is a finite cell complex in the B-rep sense. Its 0-cells are points, its edges are arcs or closed curves, and
its faces are compact surfaces, possibly with holes, whose interiors are the open faces. The underlying space $|K|$ is a
disjoint union of closed connected surfaces $|K_1|,\dots,|K_n|$ (the \emph{shells}). We write $\tau\le\sigma$ when $\tau$ is a face of $\sigma$ (including
$\tau=\sigma$). For the uniqueness theorem we assume $K$ is \emph{regular}: every closed cell $\bar\sigma$ is a
ball whose boundary is a subcomplex. B-rep complexes with seam edges, closed edges or faces with holes are not
regular, and must first be subdivided. A new cell inherits the zone of the open cell of $K$ that contains it.
This keeps the zones monotone, and a realization of the original model is a realization of the subdivision, since
$h(\inter\sigma')\subseteq h(\inter\sigma)\subseteq\Z(\sigma)$ for a new cell $\sigma'\subseteq\sigma$. The carrier
condition of Section~\ref{sec:uniqueness} must then hold for the subdivided complex. A new edge inside a face $f$
has carrier $\Z(f)$, so subdividing faces adds the hypothesis $\pi_1(\Z(f))=0$; contractible zones satisfy it. A map of $|K|$ is \emph{PL} if it is linear on the simplices of some triangulation of $|K|$ that subdivides its
cells. In dimension two, any two such triangulations are related by a cell-preserving PL homeomorphism (standard; not
separately checked against a source), so the notion does not depend on the choice. New vertices created inside an edge or face inherit its zone, so the carrier condition requires those
zones to be path-connected. The existence results say where they apply to B-rep
complexes directly.

\subsection{Geometric data and zones}
\begin{itemize}
\item Each vertex $v$ has a point $p_v$ and a tolerance $\tau_v$.
\item Each edge $e$ has a $C^2$ curve $\gamma_e\colon[0,1]\to\R^3$ and a tolerance $\tau_e$.
\item Each face $f$ has a $C^2$ surface $S_f$, a compact trimmed patch $P_f\subseteq S_f$, and a tolerance
$\tau_f$. The boundary of $P_f$ consists of the images $\beta_{f,e}=S_f\circ(\text{pcurve})$ of the pcurves of the
edges $e<f$.
\end{itemize}
The \emph{deviation} of $e$ is $\mathrm{dev}_e=\max_{f>e}\max_t|\beta_{f,e}(t)-\gamma_e(t)|$, measured at equal
parameters. OpenCascade's SameParameter property asks for $\mathrm{dev}_e\le\tau_e$ in this sense; its format
specification defines the edge tolerance as a one-sided distance from the pcurve images to the 3D
curve~\cite{OCCT}.

\begin{definition}[Zones]\label{def:zones}
We use \emph{augmented} open zones:
\begin{itemize}
\item $\Z(v)=B(p_v,\tau_v)$;
\item $\Z(e)=N_{\tau_e}(\gamma_e)\cup\bigcup_{v<e}\Z(v)$;
\item $\Z(f)=N_{\tau_f}(P_f)\cup\bigcup_{e<f}\Z(e)$.
\end{itemize}
Here $N_r$ is the open $r$-neighbourhood. Augmentation makes $\Z$ monotone: $\tau\le\sigma$ implies
$\Z(\tau)\subseteq\Z(\sigma)$.
\end{definition}

\begin{definition}[Realization]\label{def:realization}
A \emph{realization} of the tolerant B-rep $T=(K,\text{data},\tau)$ is a PL embedding $h\colon|K|\to\R^3$ with
$h(\inter\sigma)\subseteq\Z(\sigma)$ for every cell $\sigma$.
\end{definition}

\begin{remark}\label{rem:opencell}
\begin{itemize}
\item \emph{Open-cell condition.} We require only $h(\inter\sigma)\subseteq\Z(\sigma)$. The closed-cell
condition with unaugmented neighbourhoods, $h(\bar\sigma)\subseteq N_{\tau_\sigma}(\cdot)$, usually fails for every
candidate, mainly because face tolerances are tiny: $h(v)$ would have to lie within $\tau_f$ of every incident patch. With augmented zones the open-cell
condition implies $h(\bar\sigma)\subseteq\Z(\sigma)$, since $\bar\sigma$ is the union of the open cells
$\tau\le\sigma$.
\item \emph{Tameness.} Requiring $h$ to be PL rules out wild embeddings such as Alexander's horned sphere.
\end{itemize}
\end{remark}

\begin{definition}[Carriers]
The \emph{carrier} of a cell is $\C(\sigma)=\bigcup_{\tau\le\sigma}\Z(\tau)$. For augmented zones
$\C(\sigma)=\Z(\sigma)$. A continuous map $g\colon|K|\to\R^3$ is \emph{carried} if $g(\bar\sigma)\subseteq
\C(\sigma)$ for every $\sigma$. By Remark~\ref{rem:opencell}, every realization is carried.
\end{definition}

\begin{definition}[Solid]\label{def:solid}
Let $S_1,\dots,S_n\subset\R^3$ be pairwise disjoint closed surfaces. The \emph{depth} of a point $x\notin\bigcup S_j$
is the number of $j$ for which $x$ lies in the bounded component of $\R^3\setminus S_j$. The solid bounded by
the $S_j$ is the closure of the set of points of odd depth.
\end{definition}
This matches the usual B-rep convention of outer shells and cavities when the shells are consistently oriented;
a consistency check of the orientation is separate.

\section{Uniqueness for shells of positive genus}\label{sec:uniqueness}

\subsection{Hypotheses}
\begin{itemize}
\item \textbf{Carrier condition.} $\C(v)$ is path-connected for every vertex $v$, and $\pi_k(\C(\sigma))=0$ for
every $k$-cell $\sigma$ with $k\ge1$. Contractible carriers suffice.
\item \textbf{Collar condition.} For each shell $K_i$ there are a compact PL submanifold $U_i\subset\R^3$ and a
PL homeomorphism $U_i\cong M_i\times[-1,1]$ with these properties:
  \begin{itemize}
  \item $M_i\times\{0\}$ is the image of a carried PL embedding $\varphi_i\colon|K_i|\to\R^3$ (a \emph{reference
  surface});
  \item all zones of the cells of $K_i$ lie in $\inter U_i$;
  \item the $U_i$ are pairwise disjoint.
  \end{itemize}
\item \textbf{Genus.} Each $M_i$ has genus at least one.
\end{itemize}

\subsection{Lemmas}

\begin{lemma}[Carrier lemma]\label{lem:carrier}
Under the carrier condition, any two carried maps $f,g\colon|K|\to\R^3$ are homotopic by a homotopy $H$ with
$H(\bar\sigma\times I)\subseteq\C(\sigma)$ for every $\sigma$.
\end{lemma}
\begin{proof}
Define $H$ on $|K|\times\{0,1\}\cup|K^{(k)}|\times I$ by induction on $k$, keeping $H(\bar\sigma\times I)
\subseteq\C(\sigma)$, with $H(\cdot,0)=f$ and $H(\cdot,1)=g$.
\begin{itemize}
\item $k=0$: for a vertex $v$, the points $f(v),g(v)$ lie in the path-connected set $\C(v)$; join them by a
path in $\C(v)$.
\item $k\ge1$: let $\sigma$ be a $k$-cell. $\bar\sigma\times I$ is a $(k+1)$-ball whose boundary
$\bar\sigma\times\{0,1\}\cup\partial\bar\sigma\times I$ is a $k$-sphere, and $H$ is already defined on it:
  \begin{itemize}
  \item on $\bar\sigma\times\{0,1\}$ it is $f$ or $g$, with image in $\C(\sigma)$;
  \item on $\partial\bar\sigma\times I$ it is defined by induction, with image in
  $\bigcup_{\tau<\sigma}\C(\tau)\subseteq\C(\sigma)$.
  \end{itemize}
  Since $\pi_k(\C(\sigma))=0$, $H$ extends over $\bar\sigma\times I$ with image in $\C(\sigma)$.
\end{itemize}
The pieces agree on common boundaries, and a finite union of closed sets glued continuously gives a continuous
$H$.
\end{proof}

\begin{lemma}[Separation]\label{lem:separation}
Assume the carrier and collar conditions, and let $h$ be a realization. Then $G=h(|K_i|)$ is a closed surface in
$\inter U_i$, and the inclusion $G\hookrightarrow U_i$ is a homotopy equivalence.
\end{lemma}
\begin{proof}
\begin{itemize}
\item By Lemma~\ref{lem:carrier} on $K_i$, $h$ and $\varphi_i$ are homotopic through maps with image in
$\bigcup\C(\sigma)\subseteq\inter U_i$, so the homotopy takes place in $U_i$.
\item $\varphi_i\colon|K_i|\to U_i$ is a homeomorphism onto $M_i\times\{0\}$ followed by inclusion, which is a
homotopy equivalence. Hence so is $h\colon|K_i|\to U_i$.
\item Since $h$ is a homeomorphism onto $G$, the inclusion $G\hookrightarrow U_i$, which equals $h\circ h^{-1}$,
is a homotopy equivalence.
\end{itemize}
\end{proof}

\begin{lemma}[Horizontality]\label{lem:horizontal}
Under the hypotheses of Lemma~\ref{lem:separation}, if $M_i$ is not a sphere, there is an ambient isotopy of
$\R^3$, supported in $U_i$ and fixed on $\partial U_i$, that carries $G$ to $M_i\times\{0\}$.
\end{lemma}
\begin{proof}
\begin{enumerate}
\item[(a)] \emph{$G$ is incompressible.}
  \begin{itemize}
  \item By Lemma~\ref{lem:separation}, $\pi_1(G)\to\pi_1(U_i)$ is an isomorphism.
  \item The boundary of a compressing disk would be a curve that is essential in $G$ and null-homotopic in $U_i$,
  which this rules out. That is incompressibility in the sense of Waldhausen~\cite[p.~58]{Waldhausen68}.
  \item Case (b) of that definition concerns spheres and does not arise. We do not use the converse direction,
  which needs the loop theorem.
  \item $G$ is a closed orientable surface in an orientable 3-manifold, hence two-sided, and it is disjoint from
  $\partial U_i$; so it satisfies Waldhausen's conventions~\cite[\S0]{Waldhausen68}.
  \end{itemize}
\item[(b)] \emph{$G$ can be made horizontal.}
  \begin{itemize}
  \item Apply Waldhausen's Proposition~3.1~\cite[pp.~65--66]{Waldhausen68} to $U_i=M_i\times I$ ($M_i$ orientable,
  not a sphere) and the incompressible surface $G$. Since $\partial G=\emptyset$, the condition
  ``$\partial G\subseteq F\times1$'' holds vacuously.
  \item It gives an isotopy, constant on $\partial U_i$, from $G$ to a surface $G'$ whose projection
  $p\colon G'\to M_i$ is a homeomorphism onto its image.
  \item $p(G')$ is open (invariance of domain) and closed in the connected surface $M_i$, so $p(G')=M_i$.
  Hence $G'$ is the graph of a function $s\colon M_i\to(-1,1)$.
  \end{itemize}
\item[(c)] \emph{Flatten.} The PL homeomorphisms of the fibres $\{x\}\times[-1,1]$ that fix both ends and send $s(x)$
to $(1-t)s(x)$, $t\in[0,1]$, give an ambient isotopy of $U_i$, fixed on $\partial U_i$, from $G'$ to
$M_i\times\{0\}$. This isotopy is topological. Alternatively, the region between $G'$ and $M_i\times\{0\}$ is a
product region, so the two surfaces are parallel~\cite[p.~58]{Waldhausen68}.
\item[(d)] Both isotopies are fixed on $\partial U_i$ and extend by the identity to $\R^3$.
\end{enumerate}
\end{proof}

\subsection{The theorem}

\begin{theorem}[Uniqueness]\label{thm:uniqueness}
Assume the carrier condition, the collar condition, and that every shell has genus at least one. Then:
\begin{enumerate}
\item the image of any realization is ambient isotopic to $\bigcup_i M_i\times\{0\}$;
\item the images of any two realizations are ambient isotopic;
\item the solids they bound (Definition~\ref{def:solid}) are ambient isotopic.
\end{enumerate}
\end{theorem}
\begin{proof}
\begin{itemize}
\item Apply Lemma~\ref{lem:horizontal} shell by shell. The supports $U_i$ are pairwise disjoint, so the
isotopies can be run simultaneously. Their end map is a homeomorphism $\Phi$ of $\R^3$ with $\Phi(G_i)=M_i$ for
every $i$.
\item By the Jordan--Brouwer separation theorem each closed surface separates $\R^3$ into two components. $\Phi$
maps the bounded component of $\R^3\setminus G_j$ onto the bounded component of $\R^3\setminus M_j$. So depth is
preserved, and $\Phi$ maps the solid bounded by the $G_j$ onto the solid bounded by the $M_j$.
\item For two realizations with maps $\Phi$ and $\Phi'$, use $\Phi'^{-1}\circ\Phi$.
\end{itemize}
\end{proof}

\begin{remark}
\begin{itemize}
\item \emph{What is not covered.} Spheres (genus 0) are excluded, and the isotopy is of images, not of
embeddings.
\item \emph{The collar condition cannot be dropped.} The thick pipe of Section~\ref{sec:experiments} passes both
OpenCascade checks and has two realizations whose solids are not even homeomorphic.
\end{itemize}
\end{remark}

\section{The collar condition is computable}\label{sec:collar}

\subsection{Weak feature size}
Let $M\subset\R^3$ be compact, $d(p)=\dist(p,M)$, and let $\Gamma(p)$ be the set of nearest points of $p$ in $M$.
A point $p\notin M$ is \emph{critical} if $0\in\conv\{p-x:x\in\Gamma(p)\}$. The \emph{weak feature size}
$\wfs(M)$ is the infimum of the positive critical values of $d$~\cite{ChazalLieutier07}. Unlike the reach, it
stays positive at sharp edges and small fillets. A surface with a concave edge has reach $0$, and a fillet of
radius $r$ caps the reach at $r$.

\begin{lemma}[Collar]\label{lem:collar}
Let $M$ be a closed PL surface in $\R^3$, and let $0<r<\wfs(M)$.
\begin{enumerate}
\item[(a)] The closed $r$-neighbourhood $\bar N_r(M)$ is PL-homeomorphic to $M\times[-1,1]$, with $M$ going to
$M\times\{0\}$.
\item[(b)] Suppose every zone of every cell of a shell lies in $N_d(M)$ for some $d<r$, $M$ is the image of a
carried PL embedding of that shell, and the neighbourhoods $\bar N_r(M_i)$ of different shells are pairwise
disjoint. Then the collar condition of Section~\ref{sec:uniqueness} holds with $U_i=\bar N_r(M_i)$.
\end{enumerate}
\end{lemma}
\begin{proof}[Proof sketch]
\begin{itemize}
\item \emph{Small radii.} For small $\delta>0$, $\bar N_\delta(M)$ is a regular neighbourhood of $M$ relative to
$M$, by Durfee's theorem on neighbourhoods of semialgebraic sets~\cite[Props.~3.5,~3.7]{Durfee83}. A regular
neighbourhood of a closed orientable surface in $\R^3$ is a product.
\item \emph{Larger radii.} The region $\{\delta\le d\le r\}$ contains no critical point, so it is a product of
a level set with $[\delta,r]$, by the isotopy lemma of Chazal and Lieutier~\cite[Prop.~3.4]{ChazalLieutier07}.
\item \emph{Gluing.} Gluing the two products gives (a). Part (b) follows directly.
\end{itemize}
\draftnote{The application of~\cite[Prop.~3.4]{ChazalLieutier07}, which is stated for distance to the
complement of a bounded open set, to $d=\dist(\cdot,M)$ is spelled out in the repository notes (review of Lemma~1); it should be written out here. The PL (rather than
topological) form of (a) also needs a sentence.}
\end{proof}

\subsection{An explicit lower bound for polyhedral surfaces}
For $a\in M$, the \emph{conical radius} $\rho(a)$ is the supremum of the $r>0$ such that $M\cap B(a,r)=C\cap
B(a,r)$ for some closed cone $C$ with apex $a$.

\begin{lemma}[Cone lemma]\label{lem:cone}
Let $a\in M$, $r<\rho(a)$, and $p\notin M$ with $|p-a|<r/2$. Then every nearest point $x\in\Gamma(p)$ satisfies
$\langle p-x,\,p-a\rangle=|p-x|^2$. In particular $p$ is not critical.
\end{lemma}
\begin{proof}
\begin{itemize}
\item \emph{Nearest points lie in the ball.} $d(p)\le|p-a|<r/2$, so every $x\in\Gamma(p)$ has $|x-a|<r$.
\item \emph{The ray through $x$ stays in $M$.} The points $y_s=a+s(x-a)$, for $s$ near $1$, lie in the cone and
in $B(a,r)$, hence in $M$. So $s\mapsto|p-y_s|^2$ has a local minimum at $s=1$, which gives $\langle
p-x,x-a\rangle=0$.
\item \emph{Conclusion.} $\langle p-x,p-a\rangle=|p-x|^2+\langle p-x,x-a\rangle=d(p)^2>0$ for all
$x\in\Gamma(p)$. So $0$ is not in the convex hull of the vectors $p-x$.
\end{itemize}
\end{proof}

\begin{corollary}\label{cor:lambda}
Let $\lambda(M)=\inf_{x\in M}\sup_{a\in M}\bigl(\rho(a)/2-|x-a|\bigr)$. Then $d$ has no critical point $p$ with
$0<d(p)<\lambda(M)$, so $\wfs(M)\ge\lambda(M)$. If $M$ is a finite union of triangles, $\lambda(M)>0$.
\end{corollary}
\begin{proof}
\begin{itemize}
\item \emph{The bound.} Given $p$ with $d(p)<\lambda$ and $x\in\Gamma(p)$, pick $a$ with
$\rho(a)/2-|x-a|>d(p)$. Then $|p-a|<\rho(a)/2$, and Lemma~\ref{lem:cone} applies.
\item \emph{Positivity.} For a polyhedron, $\rho(a)>0$ for every $a$. The function $g(x)=\sup_a(\rho(a)/2-|x-a|)$
is $1$-Lipschitz and finite, since $\rho\le\operatorname{diam}M$, and $g(x)\ge\rho(x)/2>0$. $M$ is compact, so
$\min g>0$.
\end{itemize}
\end{proof}

Lemma~\ref{lem:cone} and the pointwise form of Corollary~\ref{cor:lambda} are machine-checked in Lean~4 with
Mathlib, for subsets of an arbitrary real inner product space (Appendix~\ref{sec:lean}).

\subsection{A certified lower bound for $\wfs$}
For a triangulated reference surface we certify ``no critical value in $(0,L)$'' with a best-first octree:
\begin{itemize}
\item \emph{Boxes near $M$} are certified by Lemma~\ref{lem:cone}, using a lower bound for $\rho(a)$ at a
candidate point $a$. The bound comes from the distance from $a$ to the nearest feature (planar region, boundary
segment, vertex) that does not contain $a$.
\item \emph{Boxes farther away} are certified by a direction test: a small linear program finds a unit vector $u$
with $\langle u,p-x\rangle>0$ for all candidate nearest points $x$ of all points $p$ of the box. The candidates
are the projections of the box centre onto the relevant features.
\item \emph{Stopping rule.} The search stops at the first box that is too small to split further and cannot be
certified. Its distance from $M$ is the certified bound $L$.
\end{itemize}
The computation is in double precision, not interval arithmetic. Table~\ref{tab:wfs} lists the results.

\begin{table}[h]
\centering
\begin{tabular}{lcc}
\toprule
Part & exact $\wfs$ & certified lower bound \\
\midrule
cube, side 10 & 5 & 4.530 \\
plate $4\times4\times1$ & 0.5 & 0.436 \\
plate $4\times4\times0.3$ & 0.15 & 0.137 \\
L-shape (wall 2) & 1 & 0.894 \\
L-shape, inner fillet 0.2 & 1 & 0.894 \\
block with slot of width 0.2 & 0.1 & 0.086 \\
cylinder $R=3$, $H=10$ & 3 & 2.512 \\
\bottomrule
\end{tabular}
\caption{Certified lower bounds for $\wfs$ of the triangulated parts. The inner fillet caps the reach at
$0.2$ but leaves the certified $\wfs$ bound unchanged.}\label{tab:wfs}
\end{table}

\section{Existence}\label{sec:existence}

We give two constructions that turn hypotheses on the data into a realization. Neither assumes that a perturbed
surface embeds as a whole. Theorem~\ref{thm:lemma0} does assume (P0) that each original patch is embedded, and two
of its hypotheses are non-local.

\subsection{Fans and cones}\label{sec:fans}

\subsubsection*{Setting}
$K$ is regular (every face a disk bounded by a single wire). For B-rep complexes:
\begin{itemize}
\item at a seam edge, E3$'$ below refers to the two sheets of the face;
\item at a closed edge, V1 counts edge ends;
\item for a face with inner wires, P0 becomes ``$Q_f$ is homeomorphic to the face''.
\end{itemize}
$\mathrm{dev}^*_e(t)$ is the deviation of the boundary points $b_{f,e}(t)=\partial P_f\cap\Pi(t)$ from $\gamma_e(t)$ in
the normal plane $\Pi(t)$; it is at most $\mathrm{dev}_e$.

\subsubsection*{Auxiliary sets}
\begin{itemize}
\item \emph{Normal disks.} For an edge $e$, choose parameters $0<a_e<b_e<1$ and a $C^1$ radius function $\rho_e(t)$.
$\rho_e$ is constant on the stubs $[t_v,a_e]$, $[b_e,t_w]$ and on a neighbourhood of $a_e$ and $b_e$; here $t_v$ is
the last parameter with $\gamma_e(t)\in B(p_v,r_v)$. $D_e(t)$ is the closed disk of radius $\rho_e(t)$ in $\Pi(t)$
centred at $\gamma_e(t)$, and $T_e=\bigcup_{t\in[a_e,b_e]}D_e(t)$.
\item \emph{Vertex regions.} For a vertex $v$, choose $r_v$. The vertex region $W_v$ is the union of the ball
$B(p_v,r_v)$ and, for each edge $e$ at $v$, the disks $D_e(t)$ for $t$ between $t_v$ and $a_e$ (or $b_e$).
\end{itemize}

\subsubsection*{Hypotheses}
E1--E3$'$ are required on $[t_v,t_w]$, the stubs included.
\begin{description}
\item[E1] $\mathrm{dev}^*_e(t)<\rho_e(t)<\min(\tau_e,\operatorname{reach}\gamma_e|_{[a_e,b_e]})$.
\item[E2] For every face $f>e$:
  \begin{enumerate}
  \item $\psi_f<\pi/2$, where $\psi_f$ is the largest angle between $\gamma_e'$ and $T P_f$ inside the tube;
  and $\rho_e k_f<1$, where $k_f=\kappa_f/\cos\psi_f$;
  \item $t\mapsto\pi_e(\beta_{f,e}(t))$ is strictly monotone. A sufficient condition is $L_e+\kappa_\gamma
  \mathrm{dev}^*_e<1$, where $L_e$ is the arc-length Lipschitz constant of $\beta_{f,e}-\gamma_e$;
  \item \emph{no return:} $\dist(\gamma_e(t),P_f\setminus C_f)>\rho_e(t)$. Here $C_f$ is the union of the section
  arcs of $P_f$ from $b_{f,e}$ to their exit points, of length at most
  $\ell_f=(\mathrm{dev}^*+\sqrt{\mathrm{dev}^{*2}+\mu\rho^2})/\mu$ with $\mu=1-\rho k_f$.
  \end{enumerate}
\item[E3$'$] The exit points $k_{f,e}(t)$ and $k_{g,e}(t)$ of the two face sections on $\partial D_e(t)$ are distinct.
A sufficient condition (E3) is $\alpha_e(t)>\varepsilon_f+\varepsilon_g$, with
$\varepsilon=\arcsin\min\{1,(\mathrm{dev}^*_e+k\ell^2/2)/\rho_e\}$. Here $\alpha_e$ is the angle between the inward
tangents of the two sections.
\item[V0] $\gamma_e$'s end at $v$ and the corners of the patches at $v$ lie in $\inter B(p_v,r_v)$.
\item[V1] (a) $W_v$ is strictly star-shaped with respect to $p_v$, and $W_v\subseteq\inter\Z(v)$.
(b) \emph{(non-local)} The construction outside $W_v$ meets $\partial W_v$ transversally in one simple closed curve
$\lambda_v$. Each incident edge contributes one point and each incident face one arc, in the cyclic order of the
link of $v$.
\item[S1] (a) The $W_v$ are pairwise disjoint. (b) The $T_e$ are pairwise disjoint. (c) $\gamma_e$ and $T_e$ miss $W_w$
for $w\notin\partial e$. (d) $P_f$ misses $T_e$ for $e\not<f$, and misses $W_v$ for $v\not<f$. (e) Outside
$\bigcup T\cup\bigcup W$, the patch pieces are pairwise disjoint. (f) $T_e\cap W_v=D_e(a_e)$ for $v<e$, and each
piece outside $W_v$ meets $W_v$ only in $\lambda_v$.
\item[P0] \emph{(non-local)} Each $P_f$ is embedded, and $P_f\cap(\bigcup T\cup\bigcup W)$ is a collar of $\partial
P_f$, so that $Q_f=P_f\setminus\inter(\bigcup T\cup\bigcup W)$ is a closed disk.
\end{description}
V1(b) and P0 are not consequences of local data. In the models of Section~\ref{sec:experiments} they were verified
only for the gapped cube, implicitly, through an exact check that the constructed surface is embedded. The
sampled checker does not test V1(b), P0, E2.3 beyond S1(d), or S1(e) for adjacent faces.

\begin{theorem}[Existence by fans and cones; conditional, proof outline]\label{thm:lemma0}
Under E1--E3$'$, V0, V1, S1 and P0, the tolerant B-rep has a realization.
\end{theorem}
\begin{proof}[Construction and proof outline]
\begin{enumerate}
\item \emph{Vertices and edges.} Put $h(v)=p_v$. On $[a_e,b_e]$, let $h(e)$ follow $\gamma_e$.
\item \emph{Fans.} In each disk $D_e(t)$, the radius from $\gamma_e(t)$ to $k_{f,e}(t)$ is assigned to $f$; these
radii form the fan strip $F_{f,e}$.
\item \emph{Faces away from the vertices.} Outside the vertex regions, the face is $Q_f\cup\bigcup_eF_{f,e}$.
\item \emph{Vertex regions.} Inside $W_v$, take the cone from $p_v$ over $\lambda_v$.
\item \emph{Injectivity.}
  \begin{itemize}
  \item The disks $D_e(t)$ are pairwise disjoint (E1 and the reach), and within one disk there are two radii meeting
  only at the centre (E3$'$).
  \item Nothing else enters the tubes or vertex regions (S1, P0, and E2.3 for $Q_f\cap F_{f,e}=k_{f,e}$).
  \item The cone is injective by star-shapedness and V1(b).
  \end{itemize}
\item \emph{Cell structure and zones.} The pieces of a face glue to a closed disk with boundary $h(\partial f)$, in
the order forced by P0 and V1(b). Each piece lies in the zone of its cell with a positive margin.
\item \emph{PL step.} The map is bi-Lipschitz: every crease has a positive opening angle, and at points where three
or more pieces meet, the tangent cone is an embedded disk. A fine secant approximation on a shape-regular
triangulation is then injective and stays in the zones.
\end{enumerate}
\draftnote{The full argument for steps 5--7 is in the repository (\texttt{notes/lemma0\_v02.md} \S\S4--5, with the
changes of \texttt{lemma0\_v03.md} \S1); it should be written out here. One gap is open: the apex estimate in the PL
step quotes a curvature $\kappa_\lambda$ of $\lambda_v$, but $\lambda_v$ is only $C^1$.}
\end{proof}

\begin{example}
On a cube whose faces are pulled apart by gaps, the construction was carried out exactly in rational
arithmetic: 216 triangles, with zone membership and embeddedness checked exactly
(\texttt{compute/lemma0\_box.py}).
\end{example}

\subsection{Re-intersection}

\begin{proposition}[Planar, degree three]\label{prop:B}
Suppose every face is a disk bounded by a single wire and is planar, $P_f\subseteq\Pi_f$, and every vertex lies on exactly three faces whose planes are
independent. Let $c_v$ be the common point of the three planes, $\ell_e=[c_v,c_w]$ for an edge $e=vw$, and
$\tilde P_f$ the polygon on the $c$'s of the vertices of $f$, in the order of $\partial f$. Assume:
\begin{itemize}
\item $c_v\in\Z(v)$, $\ell_e\subseteq\Z(e)$ and $\tilde P_f\subseteq\Z(f)$;
\item the $\tilde P_f$ are simple polygons that meet exactly in their common edges and vertices.
\end{itemize}
Then the cellular map onto $\bigcup\tilde P_f$ is a realization.
\end{proposition}
\begin{proof}
It is a PL homeomorphism onto $\bigcup\tilde P_f$, and each cell lands in its zone.
\end{proof}

\begin{proposition}[Curved, transversal]\label{prop:Bprime}
Suppose every face lies on a $C^2$ surface $S_f$, and there is a \emph{re-trimmed model}: patches $\tilde P_f\subseteq
S_f$, curves $\tilde\gamma_e$ and points $c_v$ bounding a solid, whose cell structure is isomorphic to $K$ (as a
B-rep complex, respecting seams and closed edges). Assume:
\begin{itemize}
\item the closed cells lie in their zones;
\item the cells are $C^2$ up to the boundary, the curves are regular, and there are no cusp corners;
\item the surfaces meet at a positive angle along every non-seam edge;
\item at every vertex the tangent cone is an embedded PL disk.
\end{itemize}
Then a realization exists.
\end{proposition}
\begin{proof}[Proof sketch]
\begin{itemize}
\item Take $|K|$ to be the re-trimmed surface with its cells, and $h_0$ the identity.
\item By compactness, $h_0(\bar\sigma)$ has a positive margin $m$ inside the open zone $\Z(\sigma)$.
\item $h_0$ is bi-Lipschitz: within a cell by regularity, across an edge by the positive crease angle, and at a
vertex by the tangent-cone condition.
\item A secant approximation on a fine shape-regular curved triangulation is injective once its Lipschitz
distance to $h_0$ is below the bi-Lipschitz constant, and it stays in the zones once its distance is below $m$.
\end{itemize}
\end{proof}

\begin{remark}[Checking the hypotheses]
\begin{itemize}
\item The re-trimmed model can be computed by extending the face surfaces and rebuilding the solid
(OpenCascade's \texttt{BOPAlgo\_MakerVolume}). The hypotheses are then checked on samples
(\texttt{compute/reintersect\_curved.py}), except the vertex part of the transversality hypothesis, which is not
checked. That matters for the cylinder and the pipe, whose seam vertices lie on only two surfaces. The PL step is
not carried out.
\item For the planar case, the hypotheses of Proposition~\ref{prop:B} for the sewn box and the sewn plate were also
re-checked in exact rational arithmetic, including embeddedness.
\item Results are in Section~\ref{sec:experiments}.
\end{itemize}
\end{remark}

\section{An obstruction}\label{sec:obstruction}

\subsection{Setting}
\begin{itemize}
\item $f\ne g$ are faces whose closed cells are disks bounded by single wires.
\item Their data are planar: $P_f\subseteq\Pi_f$ and $P_g\subseteq\Pi_g$, with the planes not parallel.
\item $\lambda=\Pi_f\cap\Pi_g$, and $\alpha\in(0,\pi/2]$ is the angle between the planes.
\item $n_f,n_g$ are unit normals, and $m_f$ is a unit vector in $\Pi_f$ orthogonal to $\lambda$ (likewise $m_g$).
\item $\pi_f$ is the orthogonal projection onto $\Pi_f$, and $U_f=\bigcup_{e<f}\Z(e)$, which contains the zones
of the vertices of $f$.
\item $\ell_f$ is the boundary loop of $f$: the arcs $\beta_{f,e}$ in the cyclic order of $\partial f$, with
consecutive arcs joined by the segment between their end points.
\end{itemize}

\begin{theorem}[Obstruction]\label{thm:obstruction}
Suppose there are $X\in\lambda$ and $\rho>0$ such that, with $J_f=X+[-\rho,\rho]m_f$ and $J_g=X+[-\rho,\rho]m_g$:
\begin{description}
\item[O0] $\beta_{f,e}\subseteq\Z(e)$ for every edge $e$ of $f$, and at every vertex $v$ of $f$ the end points of
the two arcs of $f$ at $v$ lie in $\Z(v)$; likewise for $g$.
\item[O1] $J_f\cap\ell_f=\emptyset$, and the winding number of $\ell_f$ in $\Pi_f$ about the points of $J_f$ is
odd; likewise for $g$.
\item[O2] $J_f\cap\cl\,\pi_f(U_f)=\emptyset$ and $J_g\cap\cl\,\pi_g(U_g)=\emptyset$.
\item[O3] $\Z(e)$ is simply connected for every edge $e$ of $f$ or $g$.
\item[O4] $\rho\sin\alpha>\tau_f\cos\alpha+\tau_g$ and $\rho\sin\alpha>\tau_g\cos\alpha+\tau_f$.
\end{description}
Then the tolerant B-rep has no realization.
\end{theorem}

\begin{proof}
Suppose $h$ is a realization, and let $\varphi=\pi_f\circ h|_{\bar f}$.

\emph{Step 1: over $J_f$, $h(\bar f)$ lies in the slab.}
\begin{itemize}
\item Let $y\in J_f$. We have $h(\bar f)\subseteq\Z(f)=N_{\tau_f}(P_f)\cup U_f$, and by O2 no point of $U_f$ lies over
$y$.
\item Since $P_f\subseteq\Pi_f$, $\pi_f^{-1}(y)\cap N_{\tau_f}(P_f)\subseteq\sigma_y:=y+(-\tau_f,\tau_f)n_f$.
\item So $h(\bar f)\cap\pi_f^{-1}(y)=h(\bar f)\cap\sigma_y$. These points are images of interior points of $f$,
because $h(\partial f)\subseteq U_f$.
\end{itemize}

\emph{Step 2: degree one.}
\begin{itemize}
\item Parametrize $\ell_f$ over $\partial f$: a small arc $A_v$ around each vertex goes to the joining segment, and
the rest of each edge goes to $\beta_{f,e}$. Shrink $A_v$ until $h(A_v)\subseteq\Z(v)$.
\item Build a homotopy in $U_f$ from $h|_{\partial f}$ to this parametrization:
  \begin{itemize}
  \item on $A_v$, the straight-line homotopy inside the convex ball $\Z(v)$;
  \item on the rest of an edge $e$, the boundary of the square is mapped into $\Z(e)$ (by O0, monotonicity, and
  $h(\bar e)\subseteq\Z(e)$), so it extends over the square by O3.
  \end{itemize}
\item By O2 the projected homotopy avoids every $y\in J_f$. So $\varphi|_{\partial f}$ and $\ell_f$ have the same
winding number about $y$, which is odd by O1. Hence the mod-2 degree of $\varphi$ at $y$ is $1$.
\item Fix a triangulation of $\bar f$ on which $h$ is linear. For $y\notin\varphi(\text{1-skeleton})$, which
excludes only finitely many points of $J_f$ once the slice in Step 3 is generic, $\#(h(\bar f)\cap\sigma_y)$ is
odd.
\end{itemize}

\emph{Step 3: a crossing arc.}
\begin{itemize}
\item Let $\Sigma$ be the plane through $X$ orthogonal to $\lambda$; it contains $n_f,m_f,n_g,m_g$.
\item Move $X$ along $\lambda$ and shrink $\rho$ slightly to generic values; O1, O2 and O4 are open conditions.
Genericity means: $\Sigma$ contains no vertex of the triangulations and meets their edges transversally, and the
lines $\langle x-X,m_f\rangle=\pm\rho$, $\langle x-X,m_g\rangle=\pm\rho$ contain no vertex.
\item In $\Sigma$, the rectangle $R_f=J_f+[-\tau_f,\tau_f]n_f$ meets $h(\bar f)$ in a compact 1-manifold $M_f$.
\item $M_f$ misses $h(\partial f)$, since $\pi_f(R_f)=J_f$ and O2 holds. It misses the long sides of $R_f$, by
Step 1. So its boundary lies on the two short sides.
\item A generic fibre $\sigma_y$ meets $M_f$ an odd number of times (Step 2). Circle components, and arcs with both
ends on one short side, meet it an even number of times.
\item Hence some arc $a_f\subseteq M_f$ joins the two short sides of $R_f$. Likewise there is an arc $a_g$ for $g$.
\end{itemize}

\emph{Step 4: the arcs meet.}
\begin{itemize}
\item In $\Sigma$, use the coordinates $u=\langle x-X,n_f\rangle$ and $w=\langle x-X,n_g\rangle$.
\item On the short sides of $R_f$, $|w|\ge\rho\sin\alpha-\tau_f\cos\alpha>\tau_g$ by O4, with opposite signs on
the two sides. So $a_f$ contains a sub-arc in the parallelogram $Q=\{|u|\le\tau_f,|w|\le\tau_g\}$ joining
$w=-\tau_g$ to $w=\tau_g$. Symmetrically, $a_g$ contains a sub-arc in $Q$ joining $u=-\tau_f$ to $u=\tau_f$.
\item By O4, $Q$ lies in both rectangles.
\item Two paths in a square joining the two pairs of opposite sides intersect
\unverified{Maehara 1984, cited for this square lemma}. So $h(\bar f)\cap h(\bar g)\ne\emptyset$ at a point that is
the image of interior points of both $f$ and $g$ (Step 1). This contradicts injectivity.
\end{itemize}
\end{proof}

\begin{remark}
\begin{itemize}
\item \emph{Which hypotheses bind.} Face tolerances are tiny (OpenCascade's default is $10^{-7}$), so O4 is mild.
The binding conditions are O1 and O2: the patches must cross away from the shadows of the edge and vertex zones.
\item \emph{Curved patches.} The proof uses planarity in Steps 1 and 2. A version for $C^2$ patches with $\tau_f$
below their reach should follow with tubular coordinates; it is not written.
\item \emph{Standard facts used without separate checking against a source:} mod-2 degree of PL maps and its relation
to the winding number, general position for PL maps and generic planes, parity of crossings of a segment with a PL
circle, homotopy invariance of the winding number, the nerve lemma (for the O3 test in the computations), and the
square lemma cited above.
\end{itemize}
\end{remark}

\section{Experiments with OpenCascade}\label{sec:experiments}

\subsection{Set-up}
\begin{itemize}
\item \emph{Software.} All experiments use OpenCascade 8.0.1 through \texttt{cadquery-ocp}. Scripts and results are
in the repository~\cite{repo}.
\item \emph{The two checks.} ``BRepCheck'' is \texttt{BRepCheck\_Analyzer(shape).IsValid()}. ``BOP check'' is
\texttt{BRepAlgoAPI\_Check(shape, true, true)}, which also tests geometry and self-interference.
\item \emph{Sewn models} imitate imported data. Every face of an exact part is moved by a random vector of length
between $g/2$ and $g$, and the faces are sewn back with \texttt{BRepBuilderAPI\_Sewing} at tolerance $4g$. Sewing
sets each edge tolerance to $1.05$ times its own deviation estimate, which is $1.048$--$1.05$ times our measured
deviation.
\item \emph{Scope.} These are a few parts in one kernel version, not a survey of industrial models.
\end{itemize}

\subsection{The checks do not decide existence or uniqueness}
Table~\ref{tab:occ} summarises.

\begin{table}[h]
\centering\small
\begin{tabular}{>{\raggedright}p{4.4cm}cc>{\raggedright\arraybackslash}p{5.0cm}}
\toprule
Model & BRepCheck & BOP check & What is true \\
\midrule
twisted triangular prism, two faces with linked boundaries ($\tau=0.081$) & valid & self-intersection & no
realization (linking number $-1$) \\
thick pipe, tolerance $1.3\times$ hole radius & valid & valid & two realizations with non-homeomorphic solids
(knot determinant $3$ vs.\ $1$) \\
sewn wedges, $30^\circ$ and $10^\circ$, $g=0.01$ & valid & valid & no realization
(Theorem~\ref{thm:obstruction}) \\
sewn plate, thickness $0.05$, $g=0.01$ & valid & invalid & realizable (Proposition~\ref{prop:B}) \\
filleted L-shape, $\tau=0.1$ & valid & self-intersection & $\tau+\mathrm{defl}<$ certified $\wfs$ bound of
$M_{PL}$ \\
long cylinder, $\tau=3.3$ & valid & valid & $\tau\ge\wfs$(nominal) (critical point exhibited) \\
\bottomrule
\end{tabular}
\caption{OpenCascade's checks against what the theory shows.}\label{tab:occ}
\end{table}

\subsubsection*{Twisted prism}
\begin{itemize}
\item Two non-adjacent triangular faces $A,B$ have boundary loops with linking number $-1$ (by crossing count, and
by the Gauss integral: $-1.000000009$). The loops are $0.4037>2\tau=0.1615$ apart, so the unions $T_A$, $T_B$ of
their edge and vertex zones are disjoint.
\item In a realization, each $h(e)$ lies in the convex capsule $C_e=N_\tau(e)$, which contains the end balls (all
tolerances are equal here). Homotope $h(e)$ rel end points to a segment inside $C_e$, and the vertices to the
nominal ones inside the balls. This gives a homotopy of $h(\partial A)$ to $\partial A$ inside $T_A$, so in the
complement of $T_B$; likewise for $B$.
\item The linking number is invariant under such homotopies, so $\mathrm{lk}(h(\partial A),h(\partial B))=-1$, and
$h(\partial A)$ must meet the disk $h(\bar B)$. That contradicts embeddedness.
\end{itemize}

\subsubsection*{Knotted pipe}
\begin{itemize}
\item One realization keeps the nominal geometry. The other replaces the inner face by a thin tube, of radius
$0.005$, around a centreline that follows a long trefoil.
\item The script checks the zones (maximum distance $0.983\le1.3$), embeddedness ($\varepsilon\kappa_{\max}=0.31$;
non-local distance $0.036>2\varepsilon$ outside a window of $0.04<\pi/\kappa_{\max}=0.0504$) and the knot
(determinant $3$).
\item The second solid is the exterior of a knotted arc in a ball, homeomorphic to the trefoil exterior, whose
fundamental group is non-abelian \unverified{standard knot theory; textbook citation to be added}. The nominal
solid is a solid torus. So the two solids are not homeomorphic.
\item This shows that the collar hypothesis of Theorem~\ref{thm:uniqueness} cannot be dropped. Here
$\wfs\le1<1.3$.
\end{itemize}

\subsubsection*{Sewn wedges}
\begin{itemize}
\item At the knife edge of the sewn wedges, the two planar patches cross along a line that lies outside the edge
tube, by $49\%$ ($30^\circ$) and about $3\times$ ($10^\circ$) of $\tau_e$.
\item The hypotheses of Theorem~\ref{thm:obstruction} hold. For the $30^\circ$ wedge the largest feasible $\rho$ at
the positions searched is $1.6\cdot10^{-3}$, while O4 needs only $3.7\cdot10^{-7}$; for the $10^\circ$ wedge it is
$1.6\cdot10^{-2}$.
\item An independent check in exact rational arithmetic confirmed all of O0--O4 for both wedges.
\item On the sewn box and plate and on the exact box and wedge, the checker finds no crossing witness, consistent
with their being realizable.
\end{itemize}

\subsection{Which construction applies}

\begin{table}[h]
\centering\small
\begin{tabular}{>{\raggedright}p{4.0cm}>{\raggedright}p{2.8cm}>{\raggedright}p{2.6cm}>{\raggedright\arraybackslash}p{4.0cm}}
\toprule
Model & fans (Thm.~\ref{thm:lemma0}) & re-intersection & margins of re-intersection (vertex / edge / face) \\
\midrule
exact box, cylinder, pipe, slot, filleted box, $10^\circ$ wedge & pass & pass where applicable & --- \\
L-shape, all edges filleted 0.5 & pass with variable radius$^*$ & not applicable & --- \\
L-shape fused, inner fillet 0.2 & fails (2 cusp vertices, 9 V1(a) failures); passes after merging coplanar
faces & not applicable & --- \\
sewn box $g=0.01$ & fails at vertices & pass & 45\% / 13\% / 4.9\% \\
sewn plate $g=0.01$ & fails at vertices & pass & 45\% / 13\% / 5.9\% \\
sewn cylinder $g=0.01$ & fails at vertices & pass (curved) & 55\% / 35\% / 5.0\% \\
sewn cylinder $g=0.05$ & fails at vertices & pass (curved) & 55\% / 11\% / 8.9\% \\
sewn pipe $g=0.01$ & fails at vertices & pass (curved) & 53\% / 16\% / 5.1\% \\
sewn wedges & fails at vertices & fails (edge outside tube) & no realization (Thm.~\ref{thm:obstruction}) \\
sewn filleted box $g=0.001$ & fails at vertices & not applicable (tangent faces) & open \\
\bottomrule
\end{tabular}
\caption{Constructions on exact and sewn models (sampled checks, double precision). $^*$Provisional: the checker
does not yet implement the constant radius near the vertices. Face margins are the smaller of the sampled checker's
value and an independent dense re-check with exact projections (\texttt{results/reintersect\_B3\_dense.txt}); for the
sewn plate only the dense re-check is available.}\label{tab:constr}
\end{table}

\begin{itemize}
\item \emph{Why the fans fail on sewn data.} The fan construction fails at the vertices of sewn models. Sewing sets
$\tau_e=1.05\,\mathrm{dev}_e$, and the tube radius must exceed the deviation; but at a corner of a box three tubes of
radius greater than $\mathrm{dev}$ do not fit disjointly into a vertex ball of radius about $\tau_e$.
\item \emph{Re-intersection.} Re-intersection avoids this, but needs transversal faces. ``Pass (curved)'' refers to
the edge part of the transversality hypothesis only; the vertex part is not checked.
\item \emph{Sewn models with tangent edges remain open.} On the sewn filleted box, OpenCascade sets each vertex
tolerance equal to the largest incident edge tolerance, the patch corners lie up to $0.85\,\tau_v$ from the vertex
point, and none of our vertex constructions fits.
\end{itemize}

\subsection{Other observations}
\begin{itemize}
\item \emph{Thickness heuristics.} Half the minimum distance between non-adjacent faces is a common estimate of the
safe tolerance. It overestimates the cylinder (5 against $\wfs=3$), and underestimates an L-shape with a $0.05$
inner fillet by a factor of about $97$.
\item \emph{Offsets.} In its default mode, \texttt{BRepOffsetAPI\_MakeOffsetShape} returned solids with wrong
volume (errors from $44\%$ to over a thousandfold, and in one case six shells instead of one) for a filleted box at
offsets within the $\wfs$-safe range, and BRepCheck accepted them.
\item \emph{Public documentation.} We did not find a check that tolerances leave the topology unchanged in the
public documentation of OpenCascade, or in third-party documentation of ACIS and Parasolid. That says only that
the documentation does not mention one.
\end{itemize}

\section{Limitations and open problems}\label{sec:open}

\begin{enumerate}
\item \textbf{Spheres.} Theorem~\ref{thm:uniqueness} excludes shells of genus $0$. Waldhausen's horizontality
argument does not apply to them. The PL Schoenflies theorem should give the sphere case under the same
hypotheses; this is not written.
\item \textbf{Embeddings.} We prove that the images of realizations are ambient isotopic, not that the embeddings
are isotopic as maps. The two differ by a self-homeomorphism of the surface.
\item \textbf{Tight vertices.} We do not know whether sewn models with tangent edges and OpenCascade's vertex
tolerances (equal to the largest incident edge tolerance) have realizations. Two directions are open: a vertex
construction that uses the edge tubes as well as the ball, or a vertex version of the obstruction.
\item \textbf{Mixed vertices.} Vertices where tangent and sharp edges meet need a construction combining fans and
re-intersection.
\item \textbf{Curved obstruction.} Theorem~\ref{thm:obstruction} is stated for planar patches.
\item \textbf{Rigour of the computations.}
  \begin{itemize}
  \item The $\wfs$ bounds and the hypothesis checks use double precision and sampling. Exceptions: the gapped-cube
  construction, and the exact re-verification of the sewn wedges and the sewn box.
  \item Interval arithmetic and a compiled implementation are needed for industrial models.
  \end{itemize}
\item \textbf{Formalization.} The carrier lemma in general (contractible carriers), Waldhausen's theorem and the
existence constructions are not formalized.
\end{enumerate}

\draftnote{Before submission: check against the originals Moise (Geometric Topology in Dimensions 2 and 3,
ch.~17, 35), Hempel (3-Manifolds, ch.~6, 10), Fu (1985), Maehara (1984), Andersson--Peters--Stewart (2000) and
Andersson--Stewart--Zidani (2006). Revise the novelty statements in Section~\ref{sec:intro} accordingly.}


\section*{Data and code availability}
All scripts, computed results and the Lean formalization are available in the public repository~\cite{repo}.

\section*{Declaration of competing interest}
The author declares no competing interests.

\section*{Declaration of generative AI and AI-assisted technologies in the research and writing process}
During this work the author used Claude (Anthropic) to draft proofs and text, write and run the experiment scripts and
the Lean formalization, and check citations against the original sources; independent AI review passes were used to
check proofs and citations. Lean~4 with Mathlib verified the formalized lemma. The author directed the research,
reviewed and edited the content, and takes full responsibility for the publication.

\appendix
\section{A curved-triangle lemma}\label{sec:curvedtri}

Appendices~\ref{sec:proof51} and~\ref{sec:proofBprime} need maps of triangles with one curved side that reproduce a
\emph{prescribed} parametrization of that side. This is what makes maps built separately on two adjacent faces agree
on their common edge. We give an explicit construction.

\begin{lemma}\label{lem:curvedtri}
Let $\hat c\colon[t_0,t_1]\to\R^2$ be a $C^2$ regular curve, $a_1=\hat c(t_0)$, $a_2=\hat c(t_1)$, and $a_3$ a point.
\begin{itemize}
\item Put $\psi(s)=\hat c(t_0+s(t_1-t_0))-((1-s)a_1+sa_2)$ for $s\in[0,1]$.
\item In barycentric coordinates $(\lambda_1,\lambda_2,\lambda_3)$ of the reference triangle $\Delta$, put
$\mu=\lambda_1+\lambda_2$ and
\[
F(\lambda)=\lambda_1a_1+\lambda_2a_2+\lambda_3a_3+\mu^2\,\psi(\lambda_2/\mu)\qquad(F=a_3\text{ at }\mu=0).
\]
\end{itemize}
Then:
\begin{enumerate}
\item[(a)] on the side $\lambda_3=0$, $F$ is $\lambda_2\mapsto\hat c(t_0+\lambda_2(t_1-t_0))$;
\item[(b)] on the sides $\lambda_1=0$ and $\lambda_2=0$, $F$ is affine;
\item[(c)] $F$ is $C^1$ on the closed triangle $\Delta$, including at its corners.
\end{enumerate}
Moreover, consider a family of such triangles in which the chord triangles $(a_1,a_2,a_3)$ are shape-regular with
diameter $h\to0$, and $|t_1-t_0|$ times the speed of $\hat c$ is $O(h)$, with uniform $C^2$ bounds on $\hat c$. Then for
small $h$:
\begin{enumerate}
\item[(d)] $F$ is a $C^1$ diffeomorphism of $\Delta$ onto the region bounded by the arc $\hat c([t_0,t_1])$ and the
segments $[a_2,a_3]$ and $[a_3,a_1]$, and $DF=A(I+O(h))$, where $A$ is the derivative of the affine part.
\end{enumerate}
\end{lemma}

\begin{proof}
(a) On $\lambda_3=0$ we have $\mu=1$, so $F=(1-\lambda_2)a_1+\lambda_2a_2+\psi(\lambda_2)=\hat c(t_0+\lambda_2(t_1-t_0))$.

(b) $\psi(0)=\psi(1)=0$. On $\lambda_2=0$ the correction is $\mu^2\psi(0)=0$. On $\lambda_1=0$ it is
$\mu^2\psi(1)=0$ (for $\mu>0$).

(c) For $\mu>0$, $F$ is $C^1$ (indeed $C^2$). Write $G=\mu^2\psi(s)$ with $s=\lambda_2/\mu$. Then
$dG=2\mu\psi(s)\,d\mu+\mu^2\psi'(s)\,ds$ with $|ds|=O(1/\mu)$, so $|dG|=O(\mu)\to0$ as $\mu\to0$. Hence $F$ is $C^1$ at
$a_3$ as well, with $DF(a_3)=A$. At the other two corners $\mu=1$, and $\psi$ is $C^2$ on $[0,1]$.

(d) \emph{Size of the correction.} $\psi''(s)=(t_1-t_0)^2\hat c''$ is $O(h^2)$, and $\psi(0)=\psi(1)=0$, so $|\psi|$ and
$|\psi'|$ are $O(h^2)$. Hence $|DG|=O(h^2)$ on $\Delta$, in coordinates of the unit reference triangle.

\emph{Injectivity.} By shape regularity $\|A^{-1}\|=O(1/h)$, so $DF=A+DG=A(I+O(h))$. For $x,y\in\Delta$ (which is
convex),
\[
|F(x)-F(y)|\ge|A(x-y)|-Ch^2|x-y|\ge(ch-Ch^2)|x-y|>0
\]
for small $h$. So $F$ is injective, and its differential is invertible.

\emph{Image.} $F$ maps $\partial\Delta$ homeomorphically onto the closed curve formed by the arc and the two segments.
That curve is simple for small $h$: the arc stays within $O(h^2)$ of its chord, and the two segments make angles with
the chord that are bounded below by shape regularity. An injective continuous map of a closed disk whose boundary goes
onto a Jordan curve has as image the closed region bounded by that curve.
\end{proof}

\emph{Remark.} The classical blending map with $\mu\,\psi(\lambda_2/\mu)$ in place of $\mu^2\psi(\lambda_2/\mu)$ is
exact on the curved side and affine on the straight sides. But it is only Lipschitz at $a_3$: its derivative there
depends on the direction of approach. The factor $\mu^2$ is what gives (c).

\emph{Standard facts used without separate checking against a source:} the Jordan curve theorem, and that an injective
continuous map of a closed disk onto a set bounded by its image boundary curve has the closed Jordan region as its
image (invariance of domain).

\section{Proof of Theorem~\ref{thm:lemma0}}\label{sec:proof51}

We use the notation and hypotheses of Section~\ref{sec:fans}. The \emph{exit point} $k_{f,e}(t)$ is the point where the
section of $P_f$ in $\Pi(t)$, followed from $b_{f,e}(t)$, leaves $D_e(t)$; E2 makes it unique. $m>0$ denotes the
minimum zone margin: $\rho_e<\tau_e$ and $W_v\subseteq\inter\Z(v)$, and by compactness all pieces built below lie at
distance at least $m$ inside the zones of their cells.

\subsection{Construction of $h$}
\begin{enumerate}
\item $h(v)=p_v$.
\item On $[a_e,b_e]$, $h(e)=\gamma_e$. In each $D_e(t)$, the radius $[\gamma_e(t),k_{f,e}(t)]$ is assigned to $f$. The fan
strip is $F_{f,e}=\bigcup_{t\in[a_e,b_e]}[\gamma_e(t),k_{f,e}(t)]$.
\item Outside the vertex regions, $h(f)=Q_f\cup\bigcup_{e<f}F_{f,e}$, where $Q_f=P_f\setminus\inter(\bigcup T\cup\bigcup W)$.
\item Inside $W_v$, take the cone from $p_v$ over $\lambda_v$: each $x\in\lambda_v$ is joined to $p_v$ by a segment. The
cone over the point of $\lambda_v$ contributed by $e$ is $h(e)$ between $v$ and $a_e$; the cone over the arc contributed
by $f$ is a sector of $h(f)$.
\end{enumerate}

\subsection{$h$ is a continuous injective cellular map into the zones}
\emph{Pieces and their supports.} $Q_f\subseteq P_f$ lies outside $\inter(\bigcup T\cup\bigcup W)$. $F_{f,e}\subseteq T_e$.
The cone sectors and edge stubs lie in $W_v$.

\emph{Injectivity inside $T_e$.}
\begin{itemize}
\item The disks $D_e(t)$ lie in the normal planes and have radius below the reach, so they are pairwise
disjoint~\cite{Federer59}; this does not depend on $\rho_e$ being constant.
\item Inside one disk, the construction consists of the centre and two radii to the distinct points $k_{f,e}(t)$,
$k_{g,e}(t)$ (E3$'$). They meet only at the centre.
\item No other piece enters $T_e$. Non-incident pieces are excluded by S1(b)--(d) and S1(f). The patches of $f>e$
enter $T_e$ only in the collar of $\partial P_f$ (P0 and E2.3), and that collar is removed.
\end{itemize}

\emph{Injectivity inside $W_v$.}
\begin{itemize}
\item The cone map $(s,x)\mapsto p_v+s(x-p_v)$, for $s\in(0,1]$ and $x\in\lambda_v$, is injective. Each ray from $p_v$
meets $\partial W_v$ once (V1(a)), and $\lambda_v$ is a simple curve on $\partial W_v$ (V1(b)).
\item No other piece enters $W_v$: S1(a), (c), (d) for non-incident pieces, and S1(f) for incident ones.
\end{itemize}

\emph{Injectivity elsewhere.} Outside $\bigcup T\cup\bigcup W$ the pieces are the $Q_f$. They are embedded (P0) and
pairwise disjoint (S1(e)).

\emph{Continuity and cell structure.}
\begin{itemize}
\item \emph{The pieces fit.} $Q_f$ and $F_{f,e}$ agree along $k_{f,e}$. The fans and $Q_f$ agree with the cone
along $\lambda_v$, which is by V1(b) exactly their trace on $\partial W_v$.
\item \emph{Order.}
  \begin{itemize}
  \item By P0, $\partial Q_f$ is one circle made of the curves $k_{f,e}$ and arcs $A_{f,v}$ on the spheres $\partial W_v$.
  \item By V1(b) there is exactly one arc per vertex of $f$, joining $k_{f,e}(a_e)$ to $k_{f,e'}(a_{e'})$.
  \item So these pieces form a 2-regular graph isomorphic to $\partial f$, whose only cycle is $\partial Q_f$; the
  cyclic order is forced.
  \end{itemize}
\item \emph{Exact intersections.}
  \begin{itemize}
  \item $Q_f\cap F_{f,e}=k_{f,e}$, by E2.3 and S1(f).
  \item $(Q_f\cup\text{fans})\cap\text{sector}_{f,v}=A_{f,v}$, by star-shapedness and S1(f).
  \item Gluing two disks along a boundary arc that is their whole intersection gives a disk. By induction, $h(f)$ is
  a closed disk with $\partial h(f)=\bigcup_{e<f}h(e)$.
  \end{itemize}
\item \emph{$h$ as a map.} Fix a homeomorphism $\partial f\to\partial h(f)$ that respects the cells, and extend it
over $f$ through a chart of the disk $h(f)$ (Alexander's trick).
  \begin{itemize}
  \item On the part away from the vertex regions, take the extension piecewise $C^1$, by straightening the $C^1$
  boundary of $Q_f$ in the parameter domain of $P_f$.
  \item On each cone sector, take it radial: the sector of $\bar f$ at $v$ is the cone over an arc of the link, and a
  point at radius $s$ over a link point $y$ goes to $p_v+s(\phi(y)-p_v)$. Here $\phi$ is a $C^1$ parametrization of
  the arc $A_{f,v}$ of $\lambda_v$.
  \end{itemize}
\end{itemize}

\emph{Zones.}
\begin{itemize}
\item $Q_f\subseteq P_f\subseteq\Z(f)$.
\item $F_{f,e}\subseteq T_e\subseteq N_{\rho}(\gamma_e)\subseteq\Z(e)\subseteq\Z(f)$.
\item The cone parts lie in $W_v\subseteq\Z(v)$, which is contained in $\Z(e)$ and $\Z(f)$ for incident $e$ and $f$.
\item $h(\inter e)\subseteq\gamma_e\cup W_v\subseteq\Z(e)$.
\item All inclusions hold with margin at least $m$.
\end{itemize}

\subsection{The PL step}
\emph{$h$ is bi-Lipschitz.} $h$ is injective on a compact set, so it suffices that it is locally bi-Lipschitz.
\begin{itemize}
\item \emph{At the creases,} the opening angle is positive.
  \begin{itemize}
  \item Along $k$, between fan and patch: the fan leaves $k$ in direction $-\hat r$, and the patch section leaves with
  $c'\cdot\hat r>0$ (E2), so the angle exceeds $90^\circ$.
  \item Along $\gamma_e$, between two fans: E3$'$ and compactness.
  \item Along $\partial W_v$, between the cone and the outside pieces: strict star-shapedness.
  \end{itemize}
\item \emph{Where three or more pieces meet} ($k_{f,e}(a_e)$, $\gamma_e(a_e)$, corners of $\partial W_v$), the tangent
cone is an embedded PL disk.
  \begin{itemize}
  \item At $k(a_e)$: the fan and $Q_f$ tangent sectors lie in the closed quarter-space
  $\{\langle\cdot,\gamma'\rangle\le0,\langle\cdot,\hat r\rangle\le0\}$. This uses $\rho_e'(a_e)=0$: $\rho_e$ is constant
  near $a_e$. The cone sectors lie inside it, with $-(x-p_v)$ in its interior. The only shared rays are $-\hat r$ and
  the side-wall tangent, each between adjacent sectors.
  \item At $\gamma(a_e)$: the fans lie in $\{\langle\cdot,\gamma'\rangle\ge0\}$ and the sectors in
  $\{\langle\cdot,\gamma'\rangle\le0\}$. They share only the rays $\hat r_f\ne\hat r_g$ (E3$'$).
  \end{itemize}
\item \emph{At the apex $p_v$,} $h$ is homogeneous of degree one on each sector. It is bi-Lipschitz there because
$\lambda_v$ is transverse to the rays from $p_v$ and $|x-p_v|$ is bounded below on $\lambda_v$.
\end{itemize}

\emph{Secant approximation.} Take a shape-regular triangulation of $|K|$ with these properties:
\begin{itemize}
\item it refines the decomposition into pieces;
\item it has the vertices of $K$ as vertices;
\item on each cone sector it is radial: its triangles there have the apex as a vertex and a chord of the link as the
opposite side;
\item $h$ is $C^1$ on each closed triangle away from the apexes.
\end{itemize}
Let $S$ be the secant map, linear on each triangle and equal to $h$ at the vertices.
\begin{itemize}
\item \emph{Away from the apexes.} $h$ is $C^1$ on each closed triangle. By the standard estimate,
$\operatorname{Lip}(S-h)\le C\,\omega(Dh,\mathrm{mesh})/\sin\theta_{\min}$, where $\omega$ is a modulus of continuity and
$\theta_{\min}$ the smallest angle. This tends to $0$ with the mesh size.
\item \emph{On a radial triangle at the apex,} with base chord $[y_0,y_1]$: write $\hat\phi$ for the linear interpolation
of $\phi$ on the chord. Then
\[
(S-h)(p_v+s(y-p_v))=s\,(\hat\phi(y)-\phi(y)),
\]
which is homogeneous of degree one.
  \begin{itemize}
  \item Its radial derivative is $\hat\phi(y)-\phi(y)$, of size at most $\omega(\phi',\mathrm{mesh})\cdot\mathrm{mesh}$.
  \item Its derivative along the chord, divided by the distance $s|y-p_v|$ to the apex, is at most
  $|\hat\phi'-\phi'|/|y-p_v|\le\omega(\phi',\mathrm{mesh})/\min|y-p_v|$.
  \item So $\operatorname{Lip}(S-h)\to0$ uniformly on the cone sectors.
  \end{itemize}
  This uses only that $\phi$ is $C^1$. It closes the gap noted in earlier versions, where the estimate was stated
  with a curvature of $\lambda_v$, which is only $C^1$.
\item \emph{Conclusion.} Since $\operatorname{Lip}(S-h)\to0$ and $|S-h|\to0$, for a fine enough mesh:
  \begin{itemize}
  \item $\operatorname{Lip}(S-h)$ is below the bi-Lipschitz constant of $h$, so $S$ is injective;
  \item $|S-h|<m$, so $S(\inter\sigma)\subseteq\Z(\sigma)$.
  \end{itemize}
  $S$ is PL, so it is a realization. $\square$
\end{itemize}

\emph{Standard facts used without separate checking against a source:} disjointness of normal disks below the
reach~\cite{Federer59}, the implicit function theorem (for the $C^1$ dependence of $k_{f,e}$ on $t$), Alexander's
trick, gluing disks along boundary arcs, and the secant-approximation estimate for $C^1$ maps on shape-regular
triangulations.

\section{Proof of Proposition~\ref{prop:Bprime}}\label{sec:proofBprime}

The idea: take the re-trimmed surface $\Sigma$ itself as the model of $|K|$ and the identity as a topological
realization, then replace it by a fine secant approximation. The steps parallel the PL step of
Appendix~\ref{sec:proof51}, but there are no fans or cones.

\subsection{Margin}
By C, the identity $\mathrm{id}\colon\Sigma\to\Sigma$, together with the cell isomorphism, maps each closed cell of $K$
onto the corresponding closed cell of $\Sigma$. By B1--B3, every closed cell lies in the zone of its cell. Closed cells
are compact and zones are open, so there is $m>0$ such that $N_m$ of every closed cell lies in its zone.

\subsection{A curved triangulation of $\Sigma$}
\begin{itemize}
\item \emph{Edges.} On each curve $\tilde\gamma_e$ fix one $C^2$ regular parametrization (say by arc length). Choose
vertices on it, including its end points.
\item \emph{Faces.} In the parameter domain of each $S_f$, the region of $\tilde P_f$ is bounded by the preimages of the
$\tilde\gamma_e$; they carry the pulled-back parametrizations and vertices. Across a seam the parameter domain is
glued periodically, so both sides of a seam use the same vertices.
\item \emph{Curved triangles.} Subdivide each region into curved triangles with at most one curved edge each. The curved
edges are pieces of the boundary curves, with their pulled-back parametrizations; the corners of the region (the
preimages of the $c_v$) are vertices; and the triangles are shape-regular in the parameter domain. This is possible for
a fine enough subdivision because the boundary curves are $C^2$ and meet at positive angles (R).
\item \emph{Maps.} Map each straight reference triangle onto its curved triangle by a transformation that reproduces the
prescribed parametrization on the curved edge and is affine on the straight edges~\cite{Lenoir86}
\unverified{regularity up to the corners and invertibility of Lenoir's transformation}. Compose with the
parametrization of $S_f$.
\end{itemize}
This gives an abstract simplicial complex $T$ and a homeomorphism $h_0\colon|T|\to\Sigma$.
\begin{itemize}
\item $T$ subdivides the cell structure of $\Sigma$, so by C it is a regular subdivision of $K$. It gives $|K|$ its PL
structure (Section~\ref{sec:setting}).
\item On each closed triangle, $h_0$ is $C^1$ up to the boundary, with injective differential.
\item On a shared edge, the two adjacent triangle maps agree:
  \begin{itemize}
  \item an interior straight edge of one face is affine in the parameter domain on both sides;
  \item an edge on $\tilde\gamma_e$ reproduces the same 3D parametrization from both faces (or from both sides of a seam).
  \end{itemize}
\end{itemize}

\subsection{$h_0$ is bi-Lipschitz}
Use the intrinsic length metric of $|T|$. $h_0$ is injective by C, and $|T|$ is compact, so it suffices that $h_0$ is
locally bi-Lipschitz.
\begin{itemize}
\item \emph{Inside a face,} including across interior edges of the subdivision: $h_0$ is a composition of a
bi-Lipschitz map onto the parameter region with the regular parametrization of $S_f$.
\item \emph{Across a seam:} both sides lie on the same surface $S_f$, and in the periodic parameter domain the two
triangles are adjacent, so the same argument applies; there is no crease.
\item \emph{Along a non-seam edge:} the two patches meet along $\tilde\gamma_e$ at a dihedral angle in $(0,\pi)$ (T), so
the tangent cone at each interior point of the edge is a union of two distinct half-planes.
\item \emph{At a vertex $c_v$:} by T the tangent cone is an embedded disk made of sectors with positive opening angles,
meeting only along common rays; by R the patch corners have positive angles.
\item Hence $h_0$ is locally bi-Lipschitz everywhere, and so bi-Lipschitz.
\end{itemize}

\subsection{Secant approximation}
Let $S$ be the secant map on $T$, refined if necessary: linear on each triangle and equal to $h_0$ at the vertices.
\begin{itemize}
\item On each triangle, $h_0$ is $C^1$ up to the boundary, and the triangles are shape-regular. The standard estimate
gives $\operatorname{Lip}(S-h_0)\le C\,\omega(Dh_0,\mathrm{mesh})/\sin\theta_{\min}$ and $|S-h_0|\le C\,\mathrm{mesh}$.
\item In the intrinsic metric, these per-triangle bounds are global.
\item For a fine enough mesh:
  \begin{itemize}
  \item $\operatorname{Lip}(S-h_0)$ is below the bi-Lipschitz constant of $h_0$, so $S$ is injective;
  \item $|S-h_0|<m$, so $S$ maps every closed cell into its zone, and in particular $S(\inter\sigma)\subseteq\Z(\sigma)$.
  \end{itemize}
\end{itemize}
$S$ is PL, so it is a realization. $\square$

\emph{Standard facts used without separate checking against a source:} existence of shape-regular curved
triangulations of a planar region bounded by $C^2$ curves meeting at positive angles; the per-triangle
secant-approximation estimate; that a map whose tangent cones are embedded disks with positive opening angles is
locally bi-Lipschitz; bi-Lipschitz equivalence of the intrinsic and extrinsic metrics on a compact polyhedron; and the
curved-triangle transformations cited above.

\section{Formal verification in Lean}\label{sec:lean}

The following are proved in Lean~4 with Mathlib (repository directory \texttt{lean/}). The build has no
\texttt{sorry}, and the results depend only on the axioms \texttt{propext}, \texttt{Classical.choice} and
\texttt{Quot.sound}.

\begin{center}\small
\begin{tabular}{p{6.6cm}p{6.4cm}}
\toprule
Lean statement & Content \\
\midrule
\texttt{inner\_eq\_zero\_of\_conical} & at a nearest point $x$, $\langle p-x,x-a\rangle=0$ (Lemma~\ref{lem:cone}) \\
\texttt{inner\_sub\_eq\_norm\_sq\_of\_conical} & $\langle p-x,p-a\rangle=|p-x|^2$ \\
\texttt{not\_isCritical\_of\_conical} & $p$ is not critical (Lemma~\ref{lem:cone}) \\
\texttt{not\_isCritical\_of\_conicalCover} & pointwise form of Corollary~\ref{cor:lambda} \\
\texttt{carrier\_mono}, \texttt{carrier\_eq\_zone\_of\_mono} & carriers are monotone, and equal the zones when the
zones are augmented \\
\texttt{carried\_of\_realization}, \texttt{image\_closedCell\_subset\_carrier} & realizations are carried \\
\texttt{carried\_straightLine}, \texttt{continuous\_straightLine} & Lemma~\ref{lem:carrier} when the carriers
are convex \\
\bottomrule
\end{tabular}
\end{center}

The statements are for an arbitrary real inner product space. Criticality is formalized as
$0\in\conv\{p-x:x\text{ nearest}\}$.

\emph{How the Lean statements correspond to the paper.}
\begin{itemize}
\item The hypothesis \texttt{LocallyConicalAt M a r} is radial closure of $M$ inside the open ball $B(a,r)$. The
paper's definition of $\rho(a)$ implies it for $r<\rho(a)$; the two are not the same statement.
\item Lean's \texttt{Carried} is pointwise on open cells. By \texttt{carrier\_mono} it agrees with the paper's
closed-cell definition.
\item Positivity of $\lambda(M)$ for polyhedra (Corollary~\ref{cor:lambda}) is proved on paper only.
\end{itemize}

\begin{thebibliography}{99}

\bibitem{APS00} L.-E. Andersson, T.~J. Peters, N.~F. Stewart, Equivalence of topological form for curvilinear
geometric objects, \emph{Int.\ J.\ Comput.\ Geom.\ Appl.} 10(6) (2000) 609--622. \unverified{bibliographic data
verified via Crossref; content not read}

\bibitem{ASZ07} L.-E. Andersson, N.~F. Stewart, M. Zidani, Error analysis for operations in solid modeling in the
presence of uncertainty, \emph{SIAM J.\ Sci.\ Comput.} 29(2) (2007) 811--826.

\bibitem{ChazalLieutier07} F. Chazal, A. Lieutier, Stability and computation of topological invariants of solids
in $\mathbb R^n$, \emph{Discrete Comput.\ Geom.} 37 (2007) 601--617.

\bibitem{dREG22} S. di Rocco, D. Eklund, O. Gäfvert, Sampling and homology via bottlenecks, arXiv:2209.01654.
\unverified{exact title}

\bibitem{Durfee83} A.~H. Durfee, Neighborhoods of algebraic sets, \emph{Trans.\ Amer.\ Math.\ Soc.} 276 (1983).
\unverified{page range}

\bibitem{Federer59} H. Federer, Curvature measures, \emph{Trans.\ Amer.\ Math.\ Soc.} 93(3) (1959) 418--491.

\bibitem{Fu85} J.~H.~G. Fu, Tubular neighborhoods in Euclidean spaces, \emph{Duke Math.\ J.} 52 (1985).
\unverified{whole entry}

\bibitem{Jackson95} D.~J. Jackson, Boundary representation modelling with local tolerances, in \emph{Proc.\ 3rd ACM
Symp.\ Solid Modeling and Applications (SMA '95)}, 1995, 247--254.

\bibitem{OCCT} Open Cascade, \emph{OpenCascade Technology documentation}, version 8.0.1 (modeling data: shape
tolerances; Boolean operations specification).

\bibitem{QiShapiro06} J. Qi, V. Shapiro, $\varepsilon$-topological formulation of tolerant solid modeling,
\emph{Computer-Aided Design} 38(4) (2006) 367--377.

\bibitem{QiShapiro06b} J. Qi, V. Shapiro, Geometric interoperability with epsilon solidity, \emph{J.\ Comput.\ Inf.\
Sci.\ Eng.} 6(3) (2006) 213--220.

\bibitem{SSP01} T. Sakkalis, G. Shen, N.~M. Patrikalakis, Topological and geometric properties of interval solid
models, \emph{Graphical Models} 63(3) (2001) 163--175.

\bibitem{Segal90} M. Segal, Using tolerances to guarantee valid polyhedral modeling results, \emph{Computer
Graphics (Proc.\ SIGGRAPH '90)} 24(4) (1990) 105--114.

\bibitem{Shapiro08} V. Shapiro, Homotopy conditions for tolerant geometric queries, in \emph{Reliable Implementation
of Real Number Algorithms}, LNCS 5045, Springer, 2008, 162--180.

\bibitem{SZ08} N.~F. Stewart, M. Zidani, Transfinite interpolation for well-definition in error analysis in solid
modelling, in \emph{Reliable Implementation of Real Number Algorithms}, LNCS 5045, Springer, 2008, 181--192.

\bibitem{Waldhausen68} F. Waldhausen, On irreducible 3-manifolds which are sufficiently large, \emph{Ann.\ of
Math.} 87 (1968) 56--88.

\bibitem{repo} G. Zhang, \texttt{tolerant-solids} (code, proofs, data), \url{https://github.com/eddiezgq/tolerant-solids}.

\end{thebibliography}

\end{document}
